\documentclass[10pt,letterpaper,sans]{moderncv}
\moderncvstyle{banking}
\moderncvcolor{blue}

\usepackage[scale=0.81]{geometry}
\nopagenumbers{}
\setlength{\hintscolumnwidth}{2.25cm}

\name{Krishnaa}{Vadivel}
\title{Electrical Engineer | Semiconductor Devices | Photonics | Embedded Systems}
\email{srikrishnaa.careers@gmail.com}
\phone[mobile]{516-853-5663}
\social[linkedin]{sri-krishnaa-ece-phd}
\extrainfo{\href{https://krishnaa423.github.io/personal-website/}{Website} \textbar{} \href{https://github.com/krishnaa423}{GitHub: krishnaa423} \textbar{} \href{https://leetcode.com/u/krishnaa42342}{LeetCode: krishnaa42342} \textbar{} \href{https://hub.docker.com/u/krishnaa42342}{Docker Hub: krishnaa42342}}

\begin{document}
\makecvtitle

\section{Summary}
\cvitem{}{Electrical engineering PhD researcher with Cornell and Yale ECE training across thin-film fabrication, semiconductor physics, integrated photonics, device simulation, and scientific computing. Experience spans cleanroom GaN device development, optical and materials characterization, light-matter modeling, and distributed HPC software for semiconductor and photonic systems.}

\section{Experience}
\cventry{2021--present}{PhD Research Assistant}{Yale University, Electrical Engineering}{}{}{\begin{itemize}%
\item Developed GaN-based power-electronics devices in cleanroom and MOCVD environments during the early PhD, including selective-area etching techniques and materials characterization using SEM, cathodoluminescence, and Raman spectroscopy.
\item Modeled VCSEL emission and related light-matter interactions involving coupled photon, phonon, and electron dynamics for semiconductor and photonics problems.
\item Built distributed HPC software for exciton-polaron and light-matter interaction modeling using first-principles methods, quantum many-body theory, and accelerator-enabled scientific computing.
\item Used agentic ML tooling to generate first-principles software inputs from journal literature and technical documentation.
\item Co-authored a 2025 \textit{Journal of the American Chemical Society} paper and delivered APS presentations in 2024, 2025, and 2026 on self-trapped exciton formation and related excited-state materials physics.
\end{itemize}}
\cventry{2019--2021}{Software and Embedded Systems Engineer}{Probe Technologies Inc. / Weatherford Wireline Services}{}{}{\begin{itemize}%
\item Parallelized CUDA software for analyzing and visualizing acoustic and nuclear well-tool data used in engineering diagnostics and field interpretation.
\item Applied dolfinx-based finite-element analysis to frequency studies of steel pipes in well environments.
\item Supported instrumentation-oriented engineering around deployed well tools, connecting field data, sensor systems, and numerical analysis.
\end{itemize}}
\cventry{2018--2019}{Photonics Technical Writer}{FindLight}{}{}{\begin{itemize}%
\item Wrote technical content on lasers, integrated optics, and optoelectronic components for engineering-facing audiences.
\end{itemize}}

\section{Selected Projects}
\cvitem{Integrated Photonics Band-Pass Filter}{Designed a Meep-based integrated photonics band-pass filter, then fabricated and tested the device with an erbium-doped fiber laser setup.}
\cvitem{Semiconductor Simulation}{Developed numerical modeling of semiconductor electrostatics, drift-diffusion transport, and pn-junction-style device behavior.}
\cvitem{VCSEL Light-Matter Modeling}{Built simulation tools for VCSEL emission and coupled photon, phonon, and electron interactions in semiconductor systems.}
\cvitem{Equivariant Scientific ML}{Developed PyTorch Geometric equivariant neural networks for material energy prediction and optical-property prediction.}

\section{Education}
\cventry{2021--present}{PhD, Electrical Engineering}{Yale University}{}{}{}
\cventry{2023}{MPhil, Electrical Engineering}{Yale University}{}{}{}
\cventry{2023}{MS, Electrical Engineering}{Yale University}{}{}{}
\cventry{2017--2018}{MEng, Electrical and Computer Engineering}{Cornell University}{}{}{}
\cventry{2013--2017}{BS, Electrical and Computer Engineering}{Cornell University}{}{}{}

\section{Technical Skills}
\cvitem{Fabrication / Characterization}{Cleanroom processing, MOCVD, SEM, cathodoluminescence, Raman spectroscopy, laser-bench optical testing}
\cvitem{Devices / Photonics}{Semiconductor physics, thin films, integrated photonics, lasers, VCSEL modeling, light-matter interaction}
\cvitem{Simulation}{Semiconductor device simulation, FDTD, finite-element analysis, first-principles and quantum many-body modeling}
\cvitem{Machine Learning}{PyTorch, PyTorch Geometric, equivariant graph neural networks, optical- and materials-property prediction}
\cvitem{Software / HPC}{Python, C, C++, CUDA, Linux, MPI, OpenMP, PETSc, SLEPc, Docker}

\end{document}
